\scanpage{55}
% archive sheet number: 54
Ainsi, les $A_\sigma$ forment un \uline{faisceau} sur l'ensemble simplicial $S_X$. \deleted{\unsure{On l'appelle} un \uline{faisceau \unsure{étagé}}} [On peut dire qu'on a un objet simplicial dans la catégorie des ensembles \deleted{\unsure{linéaires et anneaux|locaux et anneaux}} \inserted{\unsure{annelés}}]

\deleted{\unsure{sup} \unsure{fibré en} \illegible{3}}
\inserted{Ici, \unsure{tous|tous les} hom de \unsure{dégénérescence|dégénérescences} sont \unsure{des} \uline{isomorphismes}, pour les coefficients [\unsure{on a aussi} \unsure{un homomorphisme des} \unsure{simplexes non dégénérés}]}

Dans le \unsure{cas|propos} actuel, il y a lieu de \uline{négliger} tout ce qui provient des simplexes de la forme \deleted{$\unsure{\partial_i\partial_j\sigma}$, \unsure{où}} \deleted{\unsure{un simplexe non} \uline{\unsure{dégénéré}}} \deleted{[i.e. des chaines} $\partial_i\sigma$, où $\partial_i$ est un opérateur face qui n'est pas un opérateur ``de coin'', $\sigma$ \uline{\unsure{non|un}} \uline{\unsure{épuisé|épuré}} [\unsure{Cela|Ceci} revient à ne retenir que les chaines saturées \deleted{\unsure{\sigma}}, et \deleted{\unsure{parmi elles}} \unsure{des|dont}

\figure{\begin{tikzpicture}[x=1cm,y=.7cm]
\draw[gray] (0,3)--(4,3);\draw[line width=1.5pt] (0,2.9)--(0,3.1) (2,2.9)--(2,3.1) (4,2.9)--(4,3.1);
\draw[gray] (.15,2.93)..controls (.9,2.65) and (1.15,3.15)..(1.55,3.08)..controls (1.8,3.35) and (1.95,3.3)..(2,3.05);
\draw (1.94,3.55)--(2,3.48)--(2.06,3.55);\draw (1.94,3.28)--(2,3.20)--(2.06,3.28);\draw (1.94,2.9)--(2.06,3.1) (1.94,3.1)--(2.06,2.9);
\draw[gray] (0,2)--(4,2);\draw[line width=1.5pt] (0,1.86)--(0,2.14) (3.5,1.86)--(3.5,2.14);\draw (4,1.92)--(4,2.06);
\draw[gray] (1,1)--(5.5,1);\foreach \x in {1,1.6,2.2,2.8,3.4,4,4.6,5.5}{\draw[gray] (\x,.96)--(\x,1.15);}
\draw[gray] (1,-.4)--(6,-.4);\fill (1,-.4) circle (.035);\draw (1,-.5)--(1,-.18);\draw (6,-.4) circle (.035);
\end{tikzpicture}}
\diagram{\begin{tikzcd}[ampersand replacement=\&]
(x_0-\,-x_n)\arrow[d]\\{}\arrow[d]\\
(x_0-\widehat x_p-x_n)\\
x_0\quad x_1\quad x_n\\
x_0-x_n
\end{tikzcd}}

\unsure{on peut en faisant intervenir encore des} \unsure{isomorphismes}).
